\documentclass[11pt]{article}
\usepackage[T1]{fontenc}
\usepackage[utf8]{inputenc}
\usepackage{lmodern}
\usepackage[a4paper,margin=1in]{geometry}
\usepackage{microtype}
\usepackage{amsmath,amssymb,amsthm,mathtools}
\usepackage{booktabs,array}
\usepackage[dvipsnames]{xcolor}
\usepackage[colorlinks=true,linkcolor=MidnightBlue,citecolor=MidnightBlue,
            urlcolor=MidnightBlue]{hyperref}
\usepackage{enumitem}
\usepackage{setspace}
\onehalfspacing
\setlist[itemize]{topsep=4pt,itemsep=2pt,leftmargin=1.5em}
\setlist[enumerate]{topsep=4pt,itemsep=2pt,leftmargin=1.8em}

\newtheorem{theorem}{Theorem}[section]
\newtheorem{lemma}[theorem]{Lemma}
\newtheorem{proposition}[theorem]{Proposition}
\newtheorem{corollary}[theorem]{Corollary}
\theoremstyle{definition}
\newtheorem{definition}[theorem]{Definition}
\newtheorem{example}[theorem]{Example}
\theoremstyle{remark}
\newtheorem{remark}[theorem]{Remark}
\newtheorem{observation}[theorem]{Observation}

\newcommand{\DR}{\mathsf{DR}}
\newcommand{\PO}{\mathsf{PO}}
\newcommand{\PP}{\mathsf{PP}}
\newcommand{\PPI}{\mathsf{PPI}}
\newcommand{\EQ}{\mathsf{EQ}}
\newcommand{\RCC}{\{\DR,\PO,\PP,\PPI\}}
\newcommand{\RCCall}{\{\DR,\PO,\EQ,\PP,\PPI\}}
\newcommand{\cl}{\operatorname{cl}}
\newcommand{\Tp}{\operatorname{Tp}}
\newcommand{\tp}{\operatorname{tp}}
\newcommand{\comp}{\operatorname{comp}}
\newcommand{\inv}{\operatorname{inv}}
\newcommand{\TypeSafe}{\operatorname{TypeSafe}}
\newcommand{\ALCIRCC}[1]{\mathcal{ALCI}_{\mathsf{RCC#1}}}

\title{\textbf{A One-Node Counterexample to Strong Arc-Consistency}\\[0.4em]
\large Inter-witness pair-domains are necessary for $\ALCIRCC{5}$
realisability}
\author{Claude (Anthropic, Opus 4.7)\\
\small Prompted by Michael Wessel}
\date{April 2026}

\begin{document}
\maketitle

\begin{center}
\fbox{\parbox{0.92\textwidth}{%
\textbf{Disclaimer.}\;
This document was generated by Claude (Anthropic), prompted by
Michael Wessel.  It corrects a conjecture made in an earlier
counterexample paper by the same author~\cite{CounterEx2}.  The
results have not been independently verified.  We invite scrutiny
and corrections.}}
\end{center}

\begin{abstract}
In a recent two-node counterexample paper~\cite{CounterEx2} we
refuted the naive typed-patchwork conjecture of~\cite{PatchworkProbe}
and proposed a stronger propagation---\emph{strong arc-consistency}
(Def.~5.1)---with a companion conjecture (Proposition~5.2) that
strong arc-consistency is complete for realisability.  We did not
prove the companion; we noted the proof would amount to a repackaging
of the split-forest soundness theorem.  This paper refutes the
companion conjecture too.  We exhibit a \emph{one-node} TBox and
network that is strongly arc-consistent but unrealisable: two
pending existentials at the same node force witnesses whose
pairwise relations are all type-unsafe.  Strong arc-consistency
checks each existential against the existing network but does not
check pending existentials against \emph{each other}.  The minimal
fix is inter-witness pair-domains, which is exactly what the
split-forest paper's rank-$d$ validity conditions
encode~\cite{SplitForest}.  Combined with the two previous papers,
this completes a three-step cascade: every simplification of
split-forest validity we have tried has a counterexample strictly
smaller than the last, and the convergence point is split-forest
validity itself.
\end{abstract}

\section{Introduction}
\label{sec:intro}

This is the third in a cascade of papers attempting to find a simpler
propagation complete for $\ALCIRCC{5}$ realisability than the
split-forest paper's rank-$d$ validity conditions.

\begin{itemize}
  \item \textbf{Patchwork probe}~\cite{PatchworkProbe} conjectured
    (Proposition~7.1) that \emph{weak} arc-consistency
    (composition + type-safety over existing $V$) is complete.
  \item \textbf{Two-node counterexample}~\cite{CounterEx2} refuted
    Proposition~7.1 with a two-node, seven-axiom example and
    proposed \emph{strong} arc-consistency (Def.~5.1): weak plus
    existential-witness extensibility against every existing $z \in
    V$.  Companion Proposition~5.2 conjectured this stronger form is
    complete.
  \item \textbf{This paper} refutes Proposition~5.2 with a
    \emph{one-node} counterexample.  Strong arc-consistency checks
    each pending existential against existing nodes but not against
    \emph{other pending existentials at the same or different nodes}.
    The minimal next fix is inter-witness pair-domains.
\end{itemize}

The cascade converges: each simplification we tried has a strictly
smaller counterexample (7 axioms $\to$ 7 axioms, 2 nodes $\to$ 1
node), and each correction adds exactly the check the previous
definition missed.  The convergence point is the split-forest paper's
rank-$d$ validity conditions, which include pair-domains for every
pair of (existing or pending) rank-$d$ states.

Structurally: \emph{the split-forest rank-$d$ validity is not
cosmetic; it is the minimum viable propagation}.

\paragraph{What this does and does not say.}
Same caveat as the two-node paper.  This is a correction to a
specific conjecture.  It does not falsify $\ALCIRCC{5}$
decidability.  The split-forest soundness
theorem~\cite{SplitForest} already uses inter-witness pair-domains
(via rank-$d$ witness menus); only our simplified propagation
definitions miss them.

\section{Strong arc-consistency recalled}
\label{sec:strong}

\begin{definition}[Strong arc-consistency; \cite{CounterEx2}, Def.~5.1]
\label{def:strong}
A typed network $(V, D, \tau)$ is \emph{strongly arc-consistent}
if it is weakly arc-consistent (composition + type-safety
fixed-point with non-empty domains over $V$) and, in addition:
for every $x \in V$ and every $\exists R. E \in \tau(x)$, there
exists $\gamma \in \Tp(C_0, T)$ with $E \in \gamma$ and the
forward-propagated universals $\{F : \forall R. F \in \tau(x)\}
\subseteq \gamma$, such that for every $z \in V \setminus \{x\}$:
\[
  \comp(\inv(R), D(x, z)) \cap \TypeSafe(\gamma, \tau(z))
  \;\neq\; \varnothing.
\]
\end{definition}

The existential clause iterates the check over every $z \in V
\setminus \{x\}$---that is, over every \emph{existing} node other
than $x$.  It does not iterate over other \emph{pending} existentials
at $x$ itself, nor at any other node.  We now show this is a
problem.

\section{The one-node counterexample}
\label{sec:ce}

\begin{example}[Two conflicting existentials at one node]
\label{ex:one-node}
Let $T$ be the $\ALCIRCC{5}$ TBox with concept names $A, D_1, D_2$:
\begin{align*}
  \text{(T1)} \quad A    &\sqsubseteq \exists \PP. D_1, \\
  \text{(T2)} \quad A    &\sqsubseteq \exists \PP. D_2, \\
  \text{(T3)} \quad D_1  &\sqsubseteq \forall \DR.  \neg D_2, \\
  \text{(T4)} \quad D_1  &\sqsubseteq \forall \PO.  \neg D_2, \\
  \text{(T5)} \quad D_1  &\sqsubseteq \forall \PP.  \neg D_2, \\
  \text{(T6)} \quad D_1  &\sqsubseteq \forall \PPI. \neg D_2, \\
  \text{(T7)} \quad D_1  &\sqsubseteq \neg D_2.
\end{align*}
Take $C_0 = A$ and $V = \{x\}$ with $\tau(x)$ a Hintikka type for
$(C_0, T)$ containing $A$.  Necessarily $\tau(x) \supseteq
\{A, \exists \PP. D_1, \exists \PP. D_2\}$.
\end{example}

\begin{observation}[Strong arc-consistency]
\label{obs:strong}
$(V, D, \tau)$ in Example~\ref{ex:one-node} is strongly arc-consistent.
\end{observation}

\begin{proof}
$V$ has one element.  There are no non-trivial pairs, so weak
arc-consistency is vacuous.

For the existential clause, iterate over $\exists R. E \in \tau(x)$:
\begin{itemize}
  \item $\exists \PP. D_1$: we need $\gamma_1 \in \Tp(C_0, T)$ with
    $D_1 \in \gamma_1$ and $\{F : \forall \PP. F \in \tau(x)\}
    \subseteq \gamma_1$.  $\tau(x)$ contains no $\forall \PP$-universal
    (the TBox has no axioms of the form $A \sqsubseteq \forall \PP.
    \cdot$), so the forward-propagation set is empty.  A Hintikka
    type $\gamma_1$ containing $D_1$ exists (take $\gamma_1$
    saturating $\{D_1, \forall \DR. \neg D_2, \forall \PO. \neg D_2,
    \forall \PP. \neg D_2, \forall \PPI. \neg D_2, \neg D_2, \ldots\}$;
    this is internally consistent).  The quantifier over
    $z \in V \setminus \{x\}$ is vacuous ($V \setminus \{x\} =
    \varnothing$), so the condition holds.
  \item $\exists \PP. D_2$: analogous.  A Hintikka type $\gamma_2$
    containing $D_2$ exists (e.g., $\{D_2, \neg D_1, \ldots\}$
    saturating).  Condition holds.
\end{itemize}
Both existentials pass.  Strong arc-consistency holds.
\end{proof}

\begin{observation}[Unrealisability]
\label{obs:unrealisable-one}
There is no model $M \models T$ containing an element $x^M$ of
type $\tau(x)$.
\end{observation}

\begin{proof}
Suppose $M$ were such a model.  By (T1) and (T2) there exist $y_1,
y_2 \in M$ with $\PP(x, y_i)$ and $D_i \in \tp^M(y_i)$ for $i = 1,
2$.  By (T7) applied at $y_1$, $\neg D_2 \in \tp^M(y_1)$.  Hence
$y_1 \neq y_2$ (they differ in whether they contain $D_2$).

Consider $R := \rho^M(y_1, y_2) \in \RCCall$.
\begin{itemize}
  \item $R = \DR$: (T3) applied at $y_1$ (which has $\forall \DR.
    \neg D_2$) forces $\neg D_2 \in \tp^M(y_2)$; but $D_2 \in
    \tp^M(y_2)$.  Contradiction.
  \item $R = \PO$: (T4) analogous.
  \item $R = \PP$: $\PP(y_1, y_2)$ means $y_2$ is $y_1$'s
    $\PP$-successor.  (T5) at $y_1$ gives $\forall \PP. \neg D_2$,
    so $y_2$ has $\neg D_2$.  Contradiction.
  \item $R = \PPI$: $\PPI(y_1, y_2)$ means $y_1$ is $y_2$'s
    $\PP$-successor, equivalently $\PP(y_2, y_1)$, equivalently
    $y_2$ has $y_1$ as a $\PP$-successor.  But we need (T6) on
    $y_1$: $\forall \PPI. \neg D_2$.  $\PPI$-successor of $y_1$ is
    $y_2$; so $y_2$ has $\neg D_2$.  Contradiction.
  \item $R = \EQ$: $y_1 = y_2$, contradicting $y_1 \neq y_2$.
\end{itemize}
No relation is admissible.  Hence no model.
\end{proof}

\begin{corollary}
Proposition~5.2 of~\cite{CounterEx2} is false:
Example~\ref{ex:one-node} is strongly arc-consistent but
unrealisable.
\end{corollary}

\section{What went wrong}
\label{sec:analysis}

The gap is structurally similar to the one-step-up gap identified
in~\cite{CounterEx2}, but one level deeper.  Where the two-node
counterexample exposed that \emph{pending-witness vs.\ existing-node}
conflicts are not caught by weak arc-consistency, this
one-node counterexample exposes that \emph{pending-witness
vs.\ pending-witness} conflicts are not caught by strong
arc-consistency.

Strong arc-consistency checks each pending existential
$\exists R. E \in \tau(x)$ by asking: does there exist a witness
type $\gamma$ such that $\gamma$ connects to every existing node
$z \in V$?  It does not ask: does $\gamma$ connect to every
\emph{other pending witness} $\gamma'$ generated by some other
pending existential?

In the one-node example, $V = \{x\}$ has no existing nodes other
than $x$; so the existing-$z$ quantifier is vacuous and every
individually well-formed $\gamma$ passes.  The conflict is entirely
between $\gamma_1$ (witness for $\exists \PP. D_1$) and $\gamma_2$
(witness for $\exists \PP. D_2$).

\begin{remark}[Scaling pattern]
The counterexample size has decreased with each correction.
\begin{itemize}
  \item Weak $\to$ strong: the gap was between $V$ and a single
    pending witness.  The counterexample is 2 nodes, 7 axioms.
  \item Strong $\to$ stronger: the gap is between two pending
    witnesses at the same node.  The counterexample is 1 node,
    7 axioms.
  \item Stronger $\to$ even stronger would presumably be: gap between
    two pending witnesses at \emph{different} nodes, requiring a
    compositionally-induced pair-domain.  We expect this to have a
    slightly larger counterexample but the same structural flavour.
\end{itemize}
We do not attempt to construct the next-level example here; the
pattern is clear enough to motivate the correct fix directly.
\end{remark}

\section{The fix: inter-witness pair-domains}
\label{sec:fix}

The correct generalisation is to treat pending witnesses as virtual
nodes and check pairwise type-safety among them.  Concretely:

\begin{definition}[Stronger arc-consistency]
\label{def:stronger}
A typed network $(V, D, \tau)$ is \emph{stronger arc-consistent} if
there exists an assignment
\[
  \gamma : \{(x, \exists R. E) : x \in V, \; \exists R. E \in \tau(x)\}
  \;\longrightarrow\; \Tp(C_0, T)
\]
such that:
\begin{enumerate}
  \item[(i)] Each $\gamma(x, \exists R. E)$ contains $E$ and the
    forward-propagated universals $\{F : \forall R. F \in \tau(x)\}$.
  \item[(ii)] For each $(x, \exists R. E)$ and each $z \in V
    \setminus \{x\}$: $\comp(\inv(R), D(x, z)) \cap
    \TypeSafe(\gamma(x, \exists R. E), \tau(z)) \neq \varnothing$.
  \item[(iii)] For each pair of pending existentials $(x, \exists
    R_1. E_1) \neq (x', \exists R_2. E_2)$: the induced pair-domain
    $\mathrm{IPD} := \mathrm{IPD}(x, R_1, x', R_2) \cap
    \TypeSafe(\gamma(x, \exists R_1. E_1),
    \gamma(x', \exists R_2. E_2))$ is non-empty, where
    $\mathrm{IPD}(x, R_1, x', R_2)$ is the RCC5-composition-induced
    set of possible relations between the two pending witnesses
    (if $x = x'$: $\mathrm{IPD} = \comp(\inv(R_1), R_2)$; otherwise
    $\mathrm{IPD} = \comp(\comp(\inv(R_1), D(x, x')), R_2)$).
  \item[(iv)] Weak arc-consistency holds on the original $V$.
\end{enumerate}
\end{definition}

Clause (iii) is the new condition.  It says: once we commit to
witness types for each pending existential, pairs of witnesses must
be pairwise type-safe under the composition-induced domains.

\begin{observation}
Example~\ref{ex:one-node} fails stronger arc-consistency.
\end{observation}

\begin{proof}
For the two pending existentials $\exists \PP. D_1$ and $\exists
\PP. D_2$ at $x$: $\mathrm{IPD}(x, \PP, x, \PP) = \comp(\PPI, \PP)
= \RCCall$.  For every choice $\gamma_1 \ni D_1$ and $\gamma_2 \ni
D_2$: $\TypeSafe(\gamma_1, \gamma_2)$ filters as in
Observation~\ref{obs:unrealisable-one}---every atom is forbidden by
some $\forall R. \neg D_2$ universal in $\gamma_1$, and $\EQ$ is
forbidden by type-disjointness.  Hence $\mathrm{IPD} \cap
\TypeSafe = \varnothing$ for every assignment, and clause (iii)
fails.
\end{proof}

\begin{proposition}[Conjectural; equivalent to split-forest rank-$d$
validity]
\label{prop:converge}
Stronger arc-consistency, augmented with propagation to pending
witnesses of \emph{pending witnesses} (recursively up to modal
depth $d$ of the concept), is equivalent to the split-forest
paper's rank-$d$ validity conditions on an implicit
quotient~\cite{SplitForest}.
\end{proposition}

\begin{proof}[Sketch]
Unfold the witness assignment $\gamma$ recursively: witnesses of
witnesses constitute rank-$(d-1)$ states relative to $V$'s rank-$d$
states.  Clause (iii) for pending witnesses is exactly the
pair-domain condition in the split-forest paper's Definition~1.10
of valid rank-$d$ quotients.  Conditions (i)--(ii) correspond to
witness-menu relation-safety and $\Need_R$-filtering.
The correspondence is bookkeeping-exact modulo
normalisation of sibling-interface descriptors (C1)--(C5), which
are implied by clauses (iii) and their recursive extension.
\end{proof}

We do not fill in this sketch; the substance is in the split-forest
paper's soundness proof~\cite{SplitForest}.  The point of
Proposition~\ref{prop:converge} is the observation that our
three-step cascade terminates at split-forest validity---not at
some strictly simpler propagation.

\section{Implications}
\label{sec:impl}

\subsection{The cascade}

Summarising the three-paper sequence:
\begin{center}
\begin{tabular}{@{}lllll@{}}
  \toprule
  Paper & Propagation & Gap caught & Counterexample size & Conjecture \\
  \midrule
  \cite{PatchworkProbe} & weak AC & intra-$V$ type-safety & ---
    & Prop.~7.1 conj.\ (false) \\
  \cite{CounterEx2} & strong AC & + $V$-to-witness extensibility & 2 nodes, 7 ax
    & Prop.~5.2 conj.\ (false) \\
  this paper & stronger AC & + witness-to-witness pair-domains & 1 node, 7 ax
    & Prop.~\ref{prop:converge}: = split-forest validity \\
  \bottomrule
\end{tabular}
\end{center}

The pattern:
\begin{itemize}
  \item Each level checks one more layer of the witness-extension
    structure.  Each simplification that omits a layer has a
    counterexample.
  \item Counterexample \emph{size} decreases as the propagation
    strengthens, because each stronger definition catches more local
    failures, so the remaining gaps are those that exploit a purer
    version of the missing layer.
  \item The final propagation that catches everything is exactly
    split-forest rank-$d$ validity: there is no shortcut.
\end{itemize}

\subsection{For decidability}

Decidability is unaffected.  The split-forest
soundness~\cite{SplitForest} and completeness
extraction~\cite{CompletenessExtraction} together establish
decidability of $\ALCIRCC{5}$ using the full rank-$d$ validity
conditions, not any of our simplified propagations.  Our three
counterexamples refute three proposed simplifications; they do not
touch the underlying result.

\subsection{For undecidability}

The cascade is \emph{further} evidence against undecidability, in
the following sense.  Every undecidability reduction attempt we
have documented was blocked by some local consistency condition.
The three papers together exhibit the full chain of consistency
conditions that must be satisfied:

\begin{enumerate}
  \item Pairwise type-safety on $V$ (captured by weak AC).
  \item Existential-witness extensibility from $V$ (captured by
    strong AC).
  \item Inter-witness pair-domain consistency (captured by stronger
    AC).
  \item Recursive versions of (3) up to modal depth $d$ (captured
    by split-forest rank-$d$ validity).
\end{enumerate}

Each level is a binary or pair-local constraint, generated
mechanically from the TBox.  None introduces a genuinely $k$-ary
($k \geq 4$) constraint in the sense of being irreducible to the
lower levels.  The cascade concretises the ``no $k$-ary synthesis''
intuition from the probe paper~\cite{PatchworkProbe}: ALCI cannot
force arity-4 or higher constraints because every constraint the
TBox can express decomposes through the witness-extension hierarchy.

If there \emph{is} a route to undecidability for $\ALCIRCC{5}$, it
must sit outside this hierarchy.  We do not see such a route.

\subsection{For the split-forest paper}

The apparent complexity of~\cite{SplitForest}'s rank-$d$
machinery---witness menus, sibling-interface descriptors (C1)--(C5),
$\Need_R$-filtering, canonical refinements---is no longer
discretionary.  Each piece corresponds to a specific
extension-consistency check that, if omitted, admits a
counterexample of the shape documented in one of these three
papers.  The rank-$d$ hierarchy is \emph{the} minimal viable
structure for completeness.  Implementations (such as the cover-tree
tableau~\cite{CoverTree}) that do not visibly perform all these
checks are either empirically equivalent on the input distribution
tested, or silently missing cases---a possibility that the
methodological note in~\cite{CounterEx2} flagged as the
same distribution-blindness hazard that bit the pre-fix 911-test
cross-validation.

\section{Honest assessment}
\label{sec:honest}

What this paper establishes:
\begin{enumerate}
  \item A concrete one-node, seven-axiom counterexample
    (Example~\ref{ex:one-node}) that is strongly arc-consistent but
    unrealisable.  The refutation of Proposition~5.2
    of~\cite{CounterEx2} is solid.
  \item The correct fix (stronger arc-consistency,
    Definition~\ref{def:stronger}) adds an inter-witness
    pair-domain clause and is conjecturally equivalent to
    split-forest rank-$d$ validity.
  \item The three-paper cascade exhibits a pattern: every
    simplification of split-forest validity has a counterexample,
    and the counterexamples shrink as propagation strengthens.
\end{enumerate}

What we do not establish:
\begin{enumerate}
  \item A proof of Proposition~\ref{prop:converge}.  The substance
    is in~\cite{SplitForest}; we have not extracted the equivalence
    rigorously.
  \item An exhaustion of possible simplifications.  There might be
    a different simpler propagation, not in the cascade we
    considered, that catches the three counterexamples simultaneously
    without going all the way to rank-$d$ validity.  We are
    sceptical but have not proved impossibility.
  \item Any route to undecidability.  The cascade is evidence
    against, not proof against.
\end{enumerate}

\paragraph{Closing.}  Three attempts to find a simpler alternative
to split-forest validity have failed, each with a strictly smaller
counterexample than the last.  We stop here.  The next productive
question for the project is probably not ``is there an even simpler
propagation'' but ``does strong empirical validation of the
cover-tree tableau (which does \emph{not} visibly perform
inter-witness pair-domain checks) imply either (a) the cover-tree
implicitly catches these cases through some other mechanism, or
(b) the tested input distribution systematically excluded them?''
That is a question about the cover-tree implementation, not about
the split-forest semantics, and it is the right direction for the
remaining theory-vs-implementation gap documented in the overview
paper.

\begin{thebibliography}{99}

\bibitem{PatchworkProbe}
Claude (Anthropic, Opus 4.7), prompted by M.~Wessel.
\newblock Patchwork preservation under $\ALCIRCC{5}$ augmentation: A
  probe toward undecidability.
\newblock Manuscript, \texttt{papers/patchwork\_augmentation\_ALCIRCC5.pdf},
  April 2026.

\bibitem{CounterEx2}
Claude (Anthropic, Opus 4.7), prompted by M.~Wessel.
\newblock A two-node counterexample to naive typed patchwork.
\newblock Manuscript,
  \texttt{papers/typed\_patchwork\_counterexample\_ALCIRCC5.pdf},
  April 2026.

\bibitem{SplitForest}
GPT-5.4 (OpenAI), prompted by M.~Wessel.
\newblock Sibling-interface descriptors and split-forest semantics for
  $\ALCIRCC{5}$.
\newblock Manuscript,
  \texttt{papers/trees/sibling\_interface\_descriptors\_ALCIRCC5\_completed\_eqsync\_canonical\_needpatched.pdf},
  April 2026.

\bibitem{CompletenessExtraction}
Claude (Anthropic, Opus 4.6), prompted by M.~Wessel.
\newblock Completeness extraction for $\ALCIRCC{5}$.
\newblock Manuscript,
  \texttt{papers/completeness\_extraction\_ALCIRCC5.pdf}, April 2026.

\bibitem{CoverTree}
Claude (Anthropic, Opus 4.6 and Opus 4.7), prompted by M.~Wessel.
\newblock A cover-tree tableau for $\ALCIRCC{5}$.
\newblock Manuscript,
  \texttt{papers/cover\_tree\_tableau\_ALCIRCC5.pdf}, April 2026.

\end{thebibliography}

\end{document}
